\item En la sección 16.7, se dedujo la rapidez del sonido en un gas mediante el teorema impulso-momento aplicado al cilindro de gas en la figura 16.20. Ahora se encontrará la rapidez del sonido en un gas empleando un distinto enfoque basado en el elemento de gas en la figura 16.18. Proceda como sigue.
\begin{enumerate}
    \item Dibuje un diagrama de cuerpo libre para este elemento que muestre las fuerzas ejercidas sobre las superficies izquierda y derecha debido a la presión del gas sobre cualquier lado del elemento.
    \item Aplique la segunda ley de Newton al elemento para demostrar que
$$ -\frac{\partial(\Delta P)}{\partial x} A \Delta x = \rho A \Delta x \frac{\partial^2 s}{\partial t^2} $$
    \item Sustituya $\Delta P = -B \frac{\partial s}{\partial x}$ (ecuación 16.30), y obtenga la siguiente ecuación de onda para el sonido:
$$ \frac{B}{\rho} \frac{\partial^2 s}{\partial x^2} = \frac{\partial^2 s}{\partial t^2} $$
    \item Para un físico matemático, esta ecuación demuestra la existencia de ondas sonoras y determina su rapidez. Como estudiante de física, usted debe dar otro paso o dos. Sustituya en la ecuación de onda la solución de prueba $s(x,t) = s_{\text{máx}} \cos(kx-\omega t)$. Demuestre que esta función satisface la ecuación de onda, siempre que $\omega / k = v = \sqrt{B/\rho}$.
\end{enumerate}